% Original analytic teaching diagram, 2026-11-09.
\definecolor{VB}{HTML}{7188A3}\definecolor{VBF}{HTML}{E5EAF0}
\definecolor{VG}{HTML}{829C8B}\definecolor{VGF}{HTML}{E8EEE9}
\definecolor{VY}{HTML}{B4A16F}\definecolor{VYF}{HTML}{F1ECDD}
\definecolor{VR}{HTML}{A7787D}\definecolor{VRF}{HTML}{F0E5E6}
\definecolor{VP}{HTML}{9588A5}\definecolor{VPF}{HTML}{ECE8F0}
\definecolor{VI}{HTML}{3B4252}\definecolor{VM}{HTML}{6B7280}\definecolor{VL}{HTML}{4C4D4F}
\begin{tikzpicture}[x=1mm,y=1mm,>=Stealth,font=\figurefont\fontsize{8}{10}\selectfont,
lab/.style={text=VM,font=\figurefont\fontsize{7.5}{9}\selectfont},
ol/.style={text=VI,font=\figurefont\fontsize{7.5}{9}\selectfont},
grad/.style={->,draw=VP,line width=1.1pt},
good/.style={->,draw=VG,line width=1.2pt},
bd/.style={draw=VL,line width=1.5pt},
pt/.style={circle,fill=VI,inner sep=0pt,minimum size=1.8mm},
box/.style={draw=VB,fill=VBF,rounded corners=1pt,text=VI,align=center,minimum height=8mm,inner sep=2pt,line width=.7pt},
outbox/.style={draw=VL,fill=white,rounded corners=1pt,text=VI,align=center,minimum height=7mm,inner sep=2pt,line width=.6pt},
ar/.style={->,draw=VL,line width=.75pt}]
\node[box,minimum width=24mm] (a) at (0,0) {候选点可行？};
\node[box,minimum width=24mm] (b) at (0,-18) {KKT 方程有解？};
\node[box,minimum width=24mm] (c) at (0,-36) {问题是凸的？};
\node[box,minimum width=24mm] (d) at (0,-54) {临界锥上的\\Lagrange 曲率};
\node[outbox,minimum width=30mm] (a1) at (46,0) {排除};
\node[outbox,minimum width=30mm] (b1) at (46,-18) {资格成立：排除\\资格不成立：不能排除};
\node[outbox,minimum width=30mm] (c1) at (46,-36) {全局最小点};
\node[outbox,minimum width=30mm] (d1) at (46,-48) {处处为正：严格局部最小};
\node[outbox,minimum width=30mm] (d2) at (46,-58) {某方向为负且 LICQ：排除};
\node[outbox,minimum width=30mm] (d3) at (46,-68) {只有非负：需要高阶信息};
\draw[ar] (a)--node[left,lab]{是}(b); \draw[ar] (b)--node[left,lab]{有}(c); \draw[ar] (c)--node[left,lab]{否}(d);
\draw[ar] (a)--node[above,lab]{否}(a1); \draw[ar] (b)--node[above,lab]{无}(b1); \draw[ar] (c)--node[above,lab]{是}(c1);
\draw[ar] (d.east)--(d1.west); \draw[ar] (d.east)--(d2.west); \draw[ar] (d.east)--(d3.west);
\end{tikzpicture}
